\documentclass[12pt,letterpaper]{article}

% ---- Packages ----
\usepackage[utf8]{inputenc}
\usepackage[T1]{fontenc}
\usepackage{amsmath,amssymb}
\usepackage{newtxtext,newtxmath}  % Times New Roman text + math (replaces mathptmx)
\usepackage[margin=1in]{geometry}
\usepackage{setspace}
\usepackage{titlesec}
\usepackage{enumitem}
\usepackage{xr-hyper}
\usepackage{hyperref}
\externaldocument{7 - Quantum Mechanics}
\usepackage{microtype}
\usepackage{booktabs}

% ---- Section formatting ----
\titleformat{\section}{\large\bfseries\sffamily}{\thesection.}{0.5em}{}
\titleformat{\subsection}{\normalsize\bfseries\sffamily}{\thesubsection}{0.5em}{}
\titleformat{\subsubsection}{\normalsize\bfseries}{\thesubsubsection}{0.5em}{}

% ---- Spacing ----
\setstretch{1.15}
\setlength{\parskip}{0.5em}
\setlength{\parindent}{0em}

% ---- Custom commands ----
\newcommand{\rhoa}{\rho_{\!A}}
\newcommand{\Ka}{K_{\!A}}
\newcommand{\vin}{v_{\mathrm{in}}}
\newcommand{\rs}{r_s}
% \hbar is already defined in LaTeX; no redefinition needed

\begin{document}

\begin{center}
{\LARGE\bfseries Unifying Theory of Dynamic Aether Mechanics:\\
Magnetic and Electric Fields}\\[1em]
{\large Erik Otto}
\end{center}

\vspace{1em}

\begin{abstract}
Within the Dynamic Aether Mechanics framework, the electric field is identified as the
Aether flow vector---the instantaneous direction and magnitude of the medium's
displacement---while the magnetic field is the flow angle---the rotational geometry or
vorticity of that flow.\ The flow vector is sourced by the condensate's spin sector: a coherent,
net-zero-mass-flux spin-phase pump driven by the rotating winding of the bound
core ring of each charged particle,
with charge sign determined by whether the ring co-rotates (electron) or counter-
rotates (positron) with its surrounding vortex.\ This paper develops the electromagnetic
mechanism on this basis, derives the inverse-cube falloff of magnetic dipole fields,
explains ferromagnetism through Coulomb energy and Pauli exclusion, and accounts for
the Meissner effect.  Electromagnetic unification follows naturally: a stationary
charge produces a spherically symmetric wave field ($\mathbf{E}$ only), but motion
Doppler-tilts the wavefronts so that angular structure emerges ($\mathbf{B}$ appears),
with the conversion factor $v^2/c^2$ recovered from two first-order Doppler factors.
The identity $\nabla \cdot \mathbf{B} \equiv 0$ becomes a mathematical tautology
$\nabla \cdot (\nabla \times \mathbf{v}) \equiv 0$, explaining why magnetic monopoles
cannot exist.
\end{abstract}

\vspace{1em}

\bigskip
\noindent\textit{This paper is part of a series presenting the Unifying Theory of
Dynamic Aether Mechanics.  For the foundational concepts of Aether binding,
the three independent properties (density, pressure, temperature), and the
unification of force constants, see Paper~1: General Overview \cite{otto-overview}.}
\bigskip

\section{Magnetism}
\label{sec:magnetism}
% ============================================================

Magnetism in this theory arises from the \emph{angular structure} of Aether flow set in
motion by electron spin.\ Each electron acts as a vortex, spinning the surrounding Aether
into a circulatory pattern.\ While the electric field is the Aether flow vector---the
instantaneous direction and magnitude of the medium's effective displacement in the
coherent spin-phase field of the bound core ring (Section~\ref{sec:electric-field})---the
magnetic field is the flow angle: the rotational geometry of that velocity field.\ When two or more electron vortices
interact, the angular structure of their flows creates pressure differentials that produce
the forces we observe as magnetism.\ These are Bernoulli (dynamic) pressure
differentials: where opposing flows meet, kinetic energy converts to static
pressure; where flows cooperate, static pressure converts to kinetic energy.\
This mechanism operates in the incompressible-flow limit---no density change
is involved---making Bernoulli dynamics the natural description of all
electromagnetic forces in a nearly incompressible superfluid.\ In the two-fluid model
(Paper~6~\cite{otto-cosmology}), electromagnetic fields are properties of the
\emph{superfluid} component: the electron vortex is a quantized circulation in
the condensate, and the electric and magnetic fields are its velocity and
angular flow structure.\ The normal component is electromagnetically nearly decoupled in the sense
of energy exchange: direct absorption and scattering of photons by
normal-fluid quasiparticles is suppressed by
$P_{\mathrm{th}}/\Ka \sim 10^{-38}$
(where $P_{\mathrm{th}} \approx \rho_\Lambda c^2 \sim 6\times10^{-10}$~J/m$^3$
is the thermal pressure in cosmic voids and $\Ka \sim 10^{28}$~J/m$^3$;
see Paper~6~\cite{otto-cosmology}).\ This suppression ensures transparency but
does not prevent the normal fraction from modifying the background propagation
speed through the two-fluid equation of state (Paper~6, Section~5.1).


\subsection{The Basic Mechanism}

As an electron spins, it causes the Aether around it to spin as well, creating a vortex
in the surrounding medium.\ When two electrons come near each other, the Aether spinning
around them begins to interact.

When spinning in the same direction (parallel spins), the Aether between the two electrons
is impeded---the flows want to move in opposite directions at the boundary.\ This creates
a region of high pressure between them, pushing them apart.

When spinning in opposite directions (antiparallel spins), the Aether between them flows
in the same direction, allowing it to move freely and creating a region of low pressure.
The surrounding higher pressure then pushes the electrons together.

The result: alike-spinning electrons repel, and opposite-spinning electrons attract.\ This
effect is greatest at the poles of the electron, where a greater surface area of spin
interacts at higher velocities, naturally producing the strongest forces along the spin
axis.

Note that magnetism arises from the angular structure of interacting electron
flows---the pressure differentials created at the boundary between two spinning regions
of Aether.\ This is distinct from both the electric field (which arises from the
coherent spin-phase pump field of the bound core ring;
Section~\ref{sec:electric-field}) and the gravitational mechanism (which arises from
the vortex concentration and transient binding of Aether).\ All three effects originate
from the electron's interaction with Aether, but they involve geometrically distinct
aspects of the same medium: the electric field is the flow vector (the instantaneous
direction and magnitude of the medium's wave-field displacement), the magnetic field
is the flow angle (how the flow pattern rotates), and gravity is a binding phenomenon
(how each vortex concentrates and cycles Aether independently).


\subsection{Field Topology: Closed Lines and No Monopoles}

In this model, the Aether flow around each electron forms a vortex---Aether drawn inward
along the spin axis, circulated outward along the equatorial plane (or vice versa),
creating a continuous loop.\ The ``magnetic field'' is really a map of the angular
structure of this flow pattern---the vorticity of the Aether velocity field.

This immediately explains why magnetic field lines always form closed loops.\ The vortex is
a circulatory flow: Aether exits one pole, curves around through the surrounding space,
and returns at the opposite pole.\ There is no source or sink---only circulation.\ A
magnetic monopole would require a point from which Aether flow radiates outward in all
directions without returning, which is not possible for a vortex.\ The topology of
rotational flow inherently forbids monopoles.

In the flow-vector/flow-angle framework, this result is even stronger.\ The magnetic
field is the vorticity of the Aether velocity field: $\mathbf{B} \propto \nabla \times
\mathbf{v}$.\ The divergence of any curl vanishes identically---$\nabla \cdot
(\nabla \times \mathbf{v}) \equiv 0$---so $\nabla \cdot \mathbf{B} = 0$ is not a
physical law requiring explanation but a mathematical identity.\ Magnetic monopoles are
forbidden by the geometry of the medium, not by an empirical regularity.

This is consistent with the experimental fact that no magnetic monopole has ever been
observed, despite extensive searches.\ Every magnet, from a single electron to a
superconducting coil, has two poles connected by closed field lines.


\subsection{Field Falloff: The Inverse Cube Law}

The magnetic field of a single electron (a magnetic dipole) falls off as the inverse cube
of distance ($B \propto 1/r^3$), in contrast to gravity's inverse square law
($F \propto 1/r^2$).\ This difference emerges naturally from the vortex model.

Gravity is produced by the monopole effect of transient binding: each atom produces a
spherically symmetric low-pressure pulse that dissipates over the surface area of an
expanding sphere ($4\pi r^2$), giving $1/r^2$ falloff.\ The gravitational source is isotropic---it has no preferred direction.

Magnetism is produced by a dipole---a vortex with two poles and opposing flow directions.
The flow pattern has both radial and angular structure.\ At any given distance, the flows
from the two poles partially cancel (the inward flow at one pole is offset by the outward
flow at the other), and this cancellation becomes more complete with distance.\ The net
effect is an additional factor of $1/r$ beyond the geometric $1/r^2$, producing the observed
$1/r^3$ falloff.

For macroscopic magnets---where the Aether flows from enormous numbers of aligned electron
vortices combine---the field is more complex at close range but still follows dipole
falloff at distances large compared to the magnet's size.\ This is consistent with
observation.


\subsection{Amp\`ere's Law: Moving Charges and Magnetic Fields}

A moving electric charge produces a magnetic field, as described by Amp\`ere's law.\ In this
theory, the mechanism connects directly to the electric field model described below: each
charge sources a coherent spin-phase wave field (the electric field), and when the
charge moves, the emitted wavefronts are Doppler-tilted---blue-shifted ahead, red-shifted
behind---producing a net azimuthal phase gradient around the axis of motion.\ This
Doppler-induced anisotropy gives the wave field a nonzero time-averaged vorticity along
the direction of travel---creating the magnetic field observed around current-carrying
conductors.

This Doppler-tilted wave pattern produces a circular rotational structure centered on
the direction of travel---exactly the circular magnetic field observed around a
current-carrying wire.\ The faster the charge moves, the greater the Doppler
anisotropy, and the stronger the resulting magnetic field.\ For a wire carrying
current (many moving electrons), the individual circulatory effects sum, producing a net
circulatory Aether flow around the wire.

This is consistent with the right-hand rule: the orientation of the angular structure
is determined by the direction of motion relative to the flow vector.\ The superposition
of these flows produces the characteristic field pattern observed around current-carrying
conductors.


\subsection{Additivity and Material Permeation}

Magnetism is additive.\ Pressure is also additive.\ The net sum of aligned electron spin
vortices determines the total Aether flow pattern and, consequently, the total magnetic
field strength.\ When more spins align, their vortex flows reinforce, producing
proportionally stronger pressure differentials.

Magnetic fields permeate non-magnetic materials.\ Pressure does not require actual
flowing Aether to transmit force.\ A non-magnetic material would disrupt the directional
flow of Aether (with randomly oriented electrons each redirecting the flow), but the
net sum of pressure would still propagate, creating magnetic forces beyond the material.
The Aether between and around the material's atoms still transmits the pressure wave from
the external vortex, just as sound pressure passes through a wall despite being scattered
by the wall's structure.


\subsection{Ferromagnetism: Domain Alignment and Exchange Interaction}

In ferromagnetic materials (iron, nickel, cobalt), the electron spins of neighboring atoms
spontaneously align within microscopic regions called magnetic domains.\ In the Aether
model, this alignment is energetically favorable because antiparallel-spinning neighbors
create low-pressure channels between them, drawing the atoms into cooperative flow
patterns.\ Once a cluster of spins aligns, the collective vortex flow reinforces further
alignment of adjacent atoms, allowing the domain to grow.

When an external magnetic field is applied, these domains reorient to align with the
external Aether flow, and the material becomes strongly magnetic.\ When the field is
removed, the domain structure can persist---producing a permanent magnet, where the frozen
alignment of electron vortices maintains a self-sustaining Aether flow pattern.

The strength of ferromagnetic exchange coupling is naturally explained by the
combination of two mechanisms developed elsewhere in this theory.\ The exchange
interaction that drives ferromagnetic alignment is not an amplified magnetic
dipole effect---it is fundamentally a \emph{Coulomb energy difference} mediated
by the Pauli exclusion principle.

When neighboring atoms' valence electrons have parallel spins, the topological
vortex exclusion principle (Paper~7~\cite{otto-quantum}) forces them into
antisymmetric spatial configurations: same-spin electron vortices cannot overlap
in the same spatial mode, so they remain farther apart on average.\ This greater
separation reduces their mutual Coulomb repulsion (transmitted through the coherent
spin-phase field, Section~\ref{sec:electric-field}), lowering the total
energy.\ For antiparallel spins, no such topological constraint exists---the
electrons can approach more closely, increasing Coulomb repulsion.

The energy scale of this effect is set by the Coulomb interaction ($\sim$1~eV per
atom pair), not by the magnetic dipole coupling ($\sim 10^{-4}$~eV).\ This
explains the observed factor of $\sim$10{,}000 between exchange and dipole
energies without requiring any amplification mechanism---the two effects operate
through entirely different Aether processes (coherent spin-phase flux coupling
versus circulatory vortex coupling).\ The crystal-structure dependence
of ferromagnetism
follows because the exchange energy depends sensitively on orbital overlap
geometry: only certain lattice spacings and electronic configurations produce a
net energy gain from parallel alignment.


\subsection{Curie Temperature: Thermal Disruption of Alignment}

A ferromagnetic material loses its magnetism when heated above a characteristic temperature
(the Curie temperature).\ In this theory, heat is the thermal energy of Aether---the
random kinetic motion of Aether particles.\ At low temperatures, this random motion is
weak, and the organized vortex flows of aligned electron spins can maintain coherent flow
patterns (magnetic domains).

As temperature rises, the thermal energy of the surrounding Aether increases.\ The random
thermal motions of the Aether begin to disrupt the organized flow patterns, buffeting the
electron vortices and breaking their alignment.\ At the Curie temperature, the thermal
disruption is strong enough to overwhelm the cooperative flow coupling that maintains
domain structure.\ Above this temperature, the material becomes paramagnetic---the spins
can still individually respond to an external field, but they can no longer maintain
spontaneous alignment against the thermal noise.

This is directly analogous to how turbulence disrupts laminar flow in classical fluids.
An organized flow pattern (laminar flow~/~magnetic alignment) can only persist when the
driving force (vortex coupling) exceeds the disrupting force (thermal turbulence).\ The
Curie temperature marks the point of transition from ordered to disordered Aether flow.

The sharpness of this transition in many materials---where magnetism drops abruptly rather
than gradually---is consistent with a cooperative phenomenon: as some spins lose
alignment, the collective vortex flow weakens, making it easier for thermal motion to
disrupt remaining aligned spins.\ The transition cascades, producing the sharp onset
observed experimentally.


\subsection{Paramagnetism}

In paramagnetic materials, individual electron spins can partially align with an applied
magnetic field, but they cannot maintain spontaneous alignment in the absence of the field.
In the Aether model, this occurs because the thermal energy of the Aether is strong enough
to randomize the spin orientations once the organizing external flow is removed, but not
strong enough to prevent temporary alignment while the field is present.

The applied magnetic field represents an organized Aether flow that biases the orientation
of electron vortices.\ Each vortex tends to align its axis with the external flow (because
alignment reduces pressure between the vortex and the surrounding flow), but thermal
fluctuations continuously knock it out of alignment.\ The result is partial, statistical
alignment that increases with field strength and decreases with temperature---exactly as
observed.

The inverse temperature dependence of paramagnetic susceptibility (Curie's law \cite{curie1895}) follows
directly: higher temperature means more vigorous thermal disruption, less net alignment
per unit field strength, and weaker paramagnetic response.


\subsection{Diamagnetism}

Diamagnetism is the weak repulsion that all materials exhibit in a magnetic field.\ It is
the weakest form of magnetic response and is typically masked by stronger paramagnetic or
ferromagnetic effects when present.

In the Aether model, diamagnetism arises from the response of the Aether flow within
atoms when an external magnetic field (organized vorticity envelope) is applied.\ The
external vorticity phase-modulates the winding pumps of the atom's paired electrons
asymmetrically.\ Because electrons exist in paired states with opposing spins (in filled
orbitals), their coupled wave fields already produce balanced, canceling emissions.\
When an external vorticity is imposed, it advances the pump phase of the
ring aligned with it and retards the phase of the opposing ring.\ The retarded ring
loses less coherent pressure to the coupling, while the advanced ring gains more---the
net result is a small back-pressure that opposes the applied field.

This explains why diamagnetism is universal (all atoms have electron pairs), weak (it is
a second-order perturbation of already-balanced flows), and always opposing (the back-pressure always resists the external field regardless of its direction).


\subsection{The Meissner Effect: Superconductivity and Complete Flux Expulsion}

In superconductors, magnetic fields are completely expelled from the interior of the
material (the Meissner effect).\ In the Aether model, this has a clear interpretation.

At superconducting temperatures, electrons form Cooper pairs---bound pairs of electrons
with opposite spin.\ In Aether terms, each Cooper pair consists of two electron vortices
spinning in opposite directions, whose coherent wave-field emissions precisely cancel.\
The Aether within and around the pair carries, on average, zero net coherent wave
field and zero net vorticity.

When these Cooper pairs condense into a coherent superconducting state, the entire
interior of the material becomes a region of zero net coherent wave field.\ An external
magnetic field (organized vorticity envelope) cannot penetrate this region because
there are no individual, unbalanced wave sources for the external vorticity to couple
to.\ The external field is forced around the material instead, producing the observed
expulsion of magnetic flux.

The characteristic penetration depth of the magnetic field into a superconductor (the
London penetration depth) corresponds, in this model, to the distance over which the
external vorticity envelope is attenuated by interaction with the Cooper pair condensate
at the surface.\ The field doesn't stop instantaneously at the surface---it decays
exponentially over a characteristic length scale, consistent with observation.


\subsection{The Electron Magnetic Moment: Antiparallel Orientation and the \texorpdfstring{$g$}{g}-Factor}

Experiments reveal that the electron's magnetic moment is antiparallel to its spin angular
momentum---it points in the opposite direction.\ This is a quantitative constraint that
the Aether model must accommodate.

In the vortex model, the spin angular momentum of the electron corresponds to the rotation
of the electron itself (or its associated Aether structure).\ The magnetic moment, however,
is determined by the circulatory Aether flow pattern the spin produces in the surrounding
medium.\ These need not point in the same direction.\ If the electron's spin drives Aether
inward along its spin axis and outward along the equatorial plane, the resulting
circulation pattern (which determines the magnetic moment direction) would be opposite to
the spin direction---consistent with the observed antiparallel relationship.

Paper~7(a)~\cite{otto-boundcore} now derives the moment exactly: the ring's
topological winding advances azimuthally once per Compton period, giving
effective current $I = e\,m_e c^2/(2\pi\hbar)$ and moment
$\mu_M = I\,\pi r_{\mathrm{loop}}^2 = e\hbar/(2m_e) = \mu_B$, with $g = 2$
following from the half-quantum circulation---no fits, no free parameters.

The magnitude of the electron's magnetic moment is characterized by the $g$-factor.\ The
Dirac equation predicts $g = 2$ for a point particle, but the experimentally measured value
is $g = 2.00231930436118$ (known to approximately thirteen significant digits \cite{fan2023}).\ The deviation from
2---the anomalous magnetic moment---arises from virtual particle interactions
(transient binding events) that slightly modify the effective vortex flow pattern around
the electron.

This connects directly to the transient binding mechanism central to this theory.\ The
anomalous magnetic moment is calculated in quantum electrodynamics by summing the
contributions of virtual particle loops.\ Each virtual pair that forms near the electron
momentarily creates its own small vortex, slightly altering the net Aether flow pattern.
The cumulative effect of all transient binding events is a slight enhancement of the
magnetic moment beyond the bare vortex value.

The extraordinary precision with which this correction is both calculated and measured
(agreement to approximately thirteen significant digits) represents both a powerful connection
and a demanding constraint: any quantitative formalization of the Aether vortex model
must reproduce this correction.\ The fact that it arises from transient binding---the
same mechanism this theory proposes for gravity---suggests a deep connection between
magnetic and gravitational phenomena, unified through the vortex dynamics of the electron.


% ============================================================
\section{The Electric Field: A Phase-Locked Spin-Wave Force}
\label{sec:electric-field}
% ============================================================

The electric field---the force between charged particles described by the Coulomb
interaction---has a natural interpretation in the Aether framework as a
\emph{phase-locked force in the condensate's spin sector}.\ Paper~7
\cite{otto-quantum} identifies photons as spin-wave excitations of the spinor
Aether condensate; the static Coulomb interaction lives in the same sector.\ The
electric field remains identified with the \emph{Aether flow vector}, understood
as the effective velocity field that the spin sector produces through spin--orbit
coupling (see Electromagnetic Waves below): microscopically it is the
\emph{spin-phase current}, the gradient flow of the relative phase
$\varphi = \theta_\uparrow - \theta_\downarrow$ of the condensate's two
components.\ What produces it is a coherent \emph{spin-phase pump} inside the
charged particle---a rotating winding that pumps relative phase at a fixed rate,
with exactly zero time-averaged mass flux.

The channel assignment is forced, not chosen.\
Paper~7(a)~\cite{otto-boundcore} (§Channel Selection) shows that a density
(sound-wave) source of the required strength would need an oscillating energy
reservoir of order $10^{2}\,m_e c^2$: a density monopole's strength is bounded
by enclosed content, and the electron's entire budget is already spent on its
mass.\ The spin phase, by contrast, is compact---a rotating winding pumps it
indefinitely at no accumulating cost.\ The resulting interaction is the
spin-sector counterpart of the classical secondary Bjerknes force between
phase-locked sources \cite{bjerknes1906}: the same $1/r^2$ potential-theory
mathematics, transplanted from the density sector, where it cannot operate at
the observed strength, to the spin sector, where it can.

The present paper develops this mechanism \emph{generically}, using only the
requirements that each source pump the spin phase coherently at a common
clock, that it carry a discrete topological sign (the pump direction), and
that the pumps synchronize.\ All of the Coulomb-force \emph{structure}---the
$1/r^2$ falloff, the sign rule, the static field-line picture---follows from
these requirements alone and is independent of the specific topology of the
source.\ The numerical value of Coulomb's constant $k$ requires a concrete
realization: Paper~7(a)~\cite{otto-boundcore} supplies the bound core ring of
a charged lepton and obtains
$k\,e^2 = \rhoa\,\kappa_e^2\,r_{\mathrm{loop}}^2/8\pi$ (with
$\kappa_e = h/2m_A$ the circulation quantum and
$r_{\mathrm{loop}} = \hbar/m_e c$ the ring radius), reproducing the measured
value to $0.2\%$ at the reference parameter point, with the source amplitude
normalized by a single topological postulate---one winding quantum per
cycle---that simultaneously yields the Bohr magneton exactly.

Where the electric field is the flow vector of the spin-sector wave field, the
magnetic field is the \emph{flow angle}: the rotational geometry of the same
effective velocity field.\ A stationary charge produces a spherically
symmetric, static spin-current pattern with zero time-averaged vorticity:
$\mathbf{E}$ only, $\mathbf{B} = 0$.\ When that charge moves, Doppler
anisotropy of the emitted wave field imparts a net azimuthal phase gradient
around the axis of motion---angular structure appears, and the magnetic field
follows.\ The magnetic field is not a separate entity but the rotational
character of the same coherent field that constitutes the electric field.


\subsection{Required Properties of the Source}
\label{sec:ring-axial-mode}

For the phase-locked spin-sector force to reproduce the Coulomb interaction,
the charge-bearing source inside each particle must satisfy four
requirements:

\begin{enumerate}[leftmargin=2em]
  \item \textbf{Coherent spin-phase pump.}\ The source must be a monopole
    (flux) source of the condensate's relative phase: a nonzero
    $q = \oint \nabla\varphi \cdot d\mathbf{A}$ through every enclosing
    surface.\ By two-component continuity this requires the enclosed
    spin-population imbalance to be actively cycled---a pump, not a mere
    motion: translating, flexing, or rigidly rotating a static structure
    leaves all enclosed totals unchanged and produces no monopole.

  \item \textbf{Discrete topological sign invariant.}\ The source must carry
    a two-valued topological invariant that fixes the \emph{direction} of
    the pump, so that the two charge species are flux sources of opposite
    sign (e.g.\ the relative winding sense of a bound structure and its
    host vortex).

  \item \textbf{A common, universal pump clock.}\ Synchronized pumping
    requires all charges to pump at one frequency.\ The medium supplies it:
    a winding that co-rotates with its host vortex flow at the healing
    boundary turns at
    %
    \begin{equation}
      \Omega \;=\; \frac{v(\xi)}{\xi} \;=\; \frac{m_A c^2}{\hbar},
      \label{eq:pump-clock-p4}
    \end{equation}
    %
    a property of the medium alone ($v(\xi) = c/\sqrt{2}$ is exact, and the
    healing length scales inversely with $m_A c$), identical for every
    singly-wound structure regardless of its mass.\ The particle's own
    Compton frequency $\omega_C = m_e c^2/\hbar$ remains its internal
    (zitterbewegung) clock and governs the magnetic moment
    (Paper~7(a)~\cite{otto-boundcore}, ``Two Windings, Two Clocks''), but
    the force-carrying pump runs on the medium's clock---which is why
    charges of unequal mass interact without their coupling averaging
    away.

  \item \textbf{Synchronization.}\ The pumps must lock in phase.\ Each
    source is self-sustained, and mutually coupled self-sustained
    oscillators of equal frequency injection-lock at the symmetric point
    (Section~\ref{sec:self-perpetuation}); with the signs carried by the
    pump directions (Requirement~2), the locked force is proportional to
    the product of signed strengths, $F \propto q_1 q_2$.
\end{enumerate}

The bound core ring of Paper~7(a) \cite{otto-boundcore}---a quantized vortex
ring with circulation $\kappa_e = h/(2 m_A)$ whose core contains a permanent
toroidal ring of bound Aether produced by the Dynamic Casimir Effect at the
vortex boundary---satisfies all four requirements.\ Its bound core co-rotates
with the flow that created it, turning the relative-phase winding at exactly
the rate~(\ref{eq:pump-clock-p4}); the rotation sense (co- vs.\
counter-rotating, Section~\ref{sec:topological-charge}) signs the pump; and
the pump amplitude is fixed by a single topological normalization---one
winding quantum per cycle, $q_0 = 2\pi\,r_{\mathrm{loop}}$---the same
postulate that yields the Bohr magneton exactly (Paper~7(a), ``Structural
Confirmation'').

With that normalization, Paper~7(a) obtains
$k\,e^2 = \rhoa\,\kappa_e^2\,r_{\mathrm{loop}}^2/8\pi$, matching the measured
Coulomb constant to $0.2\%$ at the reference parameter point.\ We therefore
treat the numerical value of $k$ as established there and develop the
remainder of the present section in its generic form, applicable to any
topology satisfying the four requirements above.


\subsection{Topological Charge: Co-Rotation vs.\ Counter-Rotation}
\label{sec:topological-charge}

Requirement~2 of Section~\ref{sec:ring-axial-mode}---a discrete binary topological
invariant distinguishing the two charge species---is realized in the bound core
ring topology of Paper~7(a) as the \emph{relative rotation coupling} between the
ring and its surrounding vortex.\ The ring and vortex may rotate in the same sense
(co-rotating) or in opposite senses (counter-rotating).\ This relative rotation
coupling is a discrete, topological property of the bound system---there is no
continuous deformation between the two configurations---and it defines the sign of
charge:

\begin{itemize}[leftmargin=2em]
  \item \textbf{Electron} (co-rotating ring/vortex): negative charge, $q = -e$.
  \item \textbf{Positron} (counter-rotating ring/vortex): positive charge, $q = +e$.
\end{itemize}

Because the coupling is purely topological (a relative winding sign), charge has no
spatial extent.\ The spherical symmetry of the electric field at large distances is
exact, not the result of angular averaging over a distributed ring source.\ This
resolves the topological-vs-geometric charge question from
Paper~7(a)~\cite{otto-boundcore} in favor of the topological interpretation:
$F(Q^2) = 1$ exactly, consistent with electron-scattering form-factor measurements
at all tested momentum transfers.

Unlike the earlier radial-flow picture, the mechanism requires no net source or sink
of Aether at the particle---both electrons and positrons are net-zero-flow structures,
distinguished only by the relative winding of ring and vortex.

The stability of charge sign follows directly from the topological structure of the
two constituents.\ The bound core ring's winding number is fixed: it cannot change
without destroying the particle, because any continuous deformation must preserve the
topological invariant of the bound circulation.\ The surrounding host vortex's flow
direction, by contrast, can be reversed---but only through a high-energy event that
decouples the core ring from the host vortex (pair production and annihilation are
the physical realizations of such decoupling).\ In all ordinary conditions, both
components retain their formation configuration, so the sign of the ring--vortex
rotation coupling---and therefore the sign of charge---is stable over the particle's
lifetime.


\subsection{The Winding Pump: How the Ring Sources the Field}
\label{sec:piston-radiation}

In the bound core ring realization of Paper~7(a)~\cite{otto-boundcore}, the
coherent source of Requirement~1 takes a specific form.\ The bound core was
created by the host vortex flow and co-rotates with it, so the ring's
relative-phase winding rotates poloidally---about the core tube's own
axis---at the inherited rate $\Omega = v(\xi)/\xi = m_A c^2/\hbar$
(Requirement~3).\ A rigidly rotating pattern changes no enclosed totals; what
pumps is the rotation's effect at the boundary: the $m = 1$ winding rotating
at $\Omega$ makes the local relative phase at every boundary point precess
uniformly---an axisymmetric, monopole-compatible drive, equivalent to an AC
Josephson condition with bias $\hbar\Omega = m_A c^2$, the DCE bind--unbind
energy scale.\ The boundary layer therefore converts spin population
sinusoidally at $\Omega$, sloshing the enclosed imbalance and sourcing the
spin-phase flux.

The character of the pump depends only on the rotation coupling of
Section~\ref{sec:topological-charge}:

\begin{itemize}[leftmargin=2em]
  \item \textbf{Electron (co-rotating).}\ The winding rotation pumps
    relative phase with one sign: a spin-phase flux source of strength
    $-q_0$.
  \item \textbf{Positron (counter-rotating).}\ The opposite rotation sense
    pumps with the opposite sign: a source of strength $+q_0$.
\end{itemize}

The sign of charge is therefore the direction of the pump---a topological
property, not a waveform property.\ (An earlier version of this section
encoded the sign in the ``leading character'' of an emitted pressure wave;
for a monochromatic wave that character is not well defined---inverting a
sinusoid is indistinguishable from a half-period shift---whereas a pump
direction is absolute.\ The present formulation supersedes it.)\ No steady
radial flux of Aether is required; the pump transports relative phase
between the components, with zero time-averaged mass flux, consistent with
local Aether conservation.


\subsection{Self-Sustained Pumps and Synchronization}
\label{sec:self-perpetuation}

The pump does not need an external driver: the winding rotation is inherited
from the structure's formation and persists as a topologically protected
motion of a lossless condensate---an isolated charge carries its full
field.\ When two charges share the medium, each pump runs at the same
universal rate $\Omega = m_A c^2/\hbar$ (Requirement~3), because $\Omega$ is
set by the medium, not by the particle's mass.\ Mutually coupled
self-sustained oscillators of equal frequency injection-lock, and they lock
at the symmetric (in-phase) point; the quadrature configuration is the
unstable one.\ Two charges therefore settle into synchronized pumping, and
the time-averaged interaction of
Section~\ref{sec:wave-interference-force} applies with the full product of
their signed strengths.

The universality of the clock matters beyond convenience: charges of
different species (electron, muon, proton) have different internal Compton
frequencies, and species-dependent pumping rates would time-average
cross-species forces to zero.\ Because the pump clock is the medium's,
electron--muon and electron--proton Coulomb interactions are as static and
as strong as electron--electron ones, as observed.


\subsection{Attraction and Repulsion: The Flux-Source Sign Rule}
\label{sec:wave-interference-force}

The Coulomb force follows from the gradient energy of the spin-phase
field.\ Two flux sources $q_1, q_2$ superpose their $1/r^2$ spin-current
fields, and the cross term of the energy
$\tfrac{K_\varphi}{2}\!\int\!(\nabla\varphi)^2\,dV$ (with
$K_\varphi = \rhoa\hbar^2/4m_A^2$ the condensate's spin-phase stiffness;
Paper~7(a)~\cite{otto-boundcore}) evaluates to
$E_{\mathrm{int}} = K_\varphi q_1 q_2/(4\pi d)$: like signs raise the energy
as they approach (repulsion), opposite signs lower it (attraction).\ This is
the sign rule of flux sources of a gradient field---the same rule that
governs a gas of vortices---and it is opposite to the rule for passive
volume responders in a sound field (where in-phase bubbles attract): the
distinction is that a pump \emph{imposes} its flux, while a bubble
\emph{responds} to pressure.\ The classical secondary Bjerknes force
\cite{bjerknes1906} remains the mechanical ancestor of the mechanism---two
phase-locked sources coupling through a shared medium with a $1/r^2$
force---but the Aether realization transplants it to the spin sector and to
imposed-flux sources, where both the sign rule and the energetics come out
right.

With the pump direction as the charge sign, the force law for synchronized
sinusoidal pumps is $F = K_\varphi q_1 q_2/(8\pi d^{\,2})$
(Paper~7(a)~\cite{otto-boundcore} carries out the quantitative match):
electrostatics with Gauss-law charges.\ Field lines are the streamlines of
the steady spin current; they begin on sources of one sign and end on
sources of the other, and Gauss's law is the statement that the pumped flux
is conserved.


\subsection{The Electron's Three Aether Effects}
\label{sec:three-effects}

With the electric field defined as the spin-phase current and the magnetic
field as the flow angle, each electron produces three distinct Aether
phenomena, each from a different aspect of its interaction with the medium:

\begin{enumerate}[leftmargin=2em]
  \item \textbf{Spin-phase pump field} (coherent monopole sourcing of the
    condensate's relative phase by the rotating winding)---produces the
    electric field.\ Determines charge sign (via ring/vortex rotation
    coupling) and Coulomb interactions.\ Falls off as $1/r^2$ (Gauss's law
    for the conserved spin-phase flux).

  \item \textbf{Flow angle} (circulatory, from spin)---produces the intrinsic magnetic moment.\
    Determines magnetic dipole interactions with other spinning particles.\ Falls off as
    $1/r^3$ (dipole cancellation).

  \item \textbf{Vortex-boundary pair creation and transient binding}---produces gravity.\ The
    vortex core acts as a rapidly moving boundary (the DCE mechanism below the permanence
    threshold), producing the transient bind--unbind pressure wave.\ Falls off as $1/r^2$
    (spherical spreading of the monopole pressure wave), but is enormously weaker than the
    electric force for reasons developed in Section~\ref{sec:em-vs-gravity}.
\end{enumerate}

These three effects coexist because they involve distinct sectors and
geometries of the condensate: the spin-phase current (relative-phase
gradient), the flow angle (rotational geometry), and the DCE-driven binding
cycle (density sector).\ They are not in competition---a single electron
sustains all three simultaneously.


\subsection{Falloff: \texorpdfstring{$1/r^2$}{1/r2} from Gauss's Law at All Distances}
\label{sec:falloff}

In the synchronized regime the field of a charge is \emph{static}: a steady
spin-current pattern with conserved flux
$\oint \nabla\varphi \cdot d\mathbf{A} = q$ through every enclosing
surface, so $|\nabla\varphi| \propto 1/r^2$ at every distance, by Gauss's
law.\ No near-field restriction is needed, and none would suffice: the
pump's wavelength $\lambda = 2\pi c/\Omega$ is subatomic, so all macroscopic
separations lie deep in the \emph{far} field, where an unlocked oscillating
coupling would alternate in sign every half-wavelength.\ What protects the
$1/r^2$ law at macroscopic distance is the lock, not proximity: a receiver
whose pumping is exactly in phase (or anti-phase) with the locally arriving
field couples only to the $1/r^2$ component---the radiative $1/r$ component
is in quadrature and time-averages to zero.\ (An earlier version of this
section placed macroscopic Coulomb interactions in the near field; with
$\lambda$ at the subatomic scale that assignment was inverted, and the
locked-response argument above replaces it.)\ The static Gauss picture and
the locked time-average picture are the same statement in two gauges: the
Coulomb force is carried by the conserved, sourced flux, at all classical
distances.


\subsection{Coulomb vs.\ Gravity: Four Independent Differences}
\label{sec:em-vs-gravity}

The enormous strength of the Coulomb force compared to gravity (a ratio of order
$10^{42}$ for two electrons) begins with a structural difference of channel: the Coulomb source pumps
the condensate's \emph{compact} spin phase, which can be driven
indefinitely, while gravity is carried by the density channel, whose source
strength is bounded by enclosed content
(Paper~7(a)~\cite{otto-boundcore}, §Channel Selection).\ On top of this
channel asymmetry sit independent quantitative differences between the
spin-phase pump field (electric) and the bind--unbind residual pressure
wave (gravity):

\begin{enumerate}[leftmargin=2em]
  \item \textbf{Amplitude.}\ The Coulomb pump runs at its full topological
    strength (one winding quantum per cycle, $q_0 = 2\pi r_{\mathrm{loop}}$).\ The gravity wave amplitude is the residual
    asymmetry of the transient bind--unbind cycle at the vortex core boundary---a small
    fraction of the oscillation-scale amplitude.

  \item \textbf{Frequency (no longer a hierarchy factor).}\ In the winding-pump
    picture both the Coulomb pump and the gravitational bind--unbind cycle run at
    the boundary scale ($\Omega \sim m_A c^2/\hbar$, set by $v_{\mathrm{perm}}$
    and the healing length~$\xi$); the frequency ratio of earlier drafts drops
    out, and its explanatory burden moves to the channel asymmetry above.

  \item \textbf{Asymmetry.}\ The spin-phase pumping is a symmetric sinusoid whose
    positive and negative half-cycles carry full, coherent, phase-locked force.\ The
    gravitational emission is asymmetric---the residual after nearly complete bind--unbind
    cancellation---and carries only the small unrecycled component.

  \item \textbf{One-way vs.\ two-way coupling.}\ Two charges \emph{mutually} drive each
    other's pumping: each pump's phase is locked to the arriving field, creating
    the synchronized coupled-oscillator regime of
    Section~\ref{sec:self-perpetuation} that underlies the Coulomb force.\ Gravity, by
    contrast, may be below the receiver's response threshold---the bind--unbind residual
    wave amplitude may be too small to drive any resonant back-reaction in a receiver
    vortex.\ The interaction is effectively one-way (asymmetric-emission only) rather
    than a mutual coupled-oscillator.\ This difference in coupling \emph{character}, not
    just in amplitude, is a significant contributor to the observed force hierarchy.
\end{enumerate}

The $10^{42}$ ratio is the product of the channel asymmetry and these
factors.\ A first-principles determination of each---pump strength, asymmetry
fraction $\epsilon$, and response-threshold criterion---is left as an open
question (Section~\ref{sec:em-open-questions}).


\subsection{Electromagnetic Unification: How Movement Connects \texorpdfstring{$\mathbf{E}$}{E} and \texorpdfstring{$\mathbf{B}$}{B}}

The relationship between electric and magnetic fields---the foundation of Maxwell's
equations---emerges naturally from the Aether model once the electric field is
identified as the flow vector of the coherent spin-sector wave field and the
magnetic field as its angular structure.

When a charge is stationary, its spin-phase field is spherically symmetric
about the ring center.\ The instantaneous flow vector points radially outward or inward
at each phase of the oscillation, but the time-averaged vorticity of this pattern is
zero: there is no magnetic field (beyond the intrinsic spin moment of the particle
itself).

When a charge moves with velocity $\mathbf{v}$, the emitted wavefronts are
Doppler-shifted---blue-shifted ahead of the motion, red-shifted behind.\ This
kinematic anisotropy of the wave field produces a net azimuthal phase gradient
around the axis of motion, so that the time-averaged pattern acquires nonzero
vorticity along~$\mathbf{v}$.\ The coherent wave field of a moving charge has
$\nabla \times \langle \mathbf{E} \rangle \neq 0$---this is the magnetic field
produced by a moving charge, expressed as a purely kinematic consequence of the
Doppler structure of the radiated wave.\ The Doppler shift is first order in $v/c$,
and the magnetic interaction between two moving charges involves a Doppler factor
from each, recovering the observed $F_{\mathrm{mag}}/F_{\mathrm{elec}} = v_1 v_2/c^2$
(Section~\ref{sec:em-math-summary}).

This mechanism provides the physical basis for several fundamental electromagnetic
relationships:

\begin{itemize}[leftmargin=2em]
  \item \textbf{Amp\`ere's law:} A current-carrying wire (many co-moving charges) produces a
    circular magnetic field around it.\ In the Aether model, each moving charge's
    Doppler-anisotropic wave field carries a nonzero azimuthal vorticity component, and
    these individual contributions sum coherently around the wire axis to produce the
    observed circular field.

  \item \textbf{The right-hand rule:} The orientation of the Doppler-induced vorticity is
    determined by the direction of motion relative to the wave field---the geometry of
    wavefront tilting dictates the resulting flow angle.

  \item \textbf{Relativity of $\mathbf{E}$ and $\mathbf{B}$ fields:} In standard physics, a pure
    electric field in one reference frame appears as a combination of electric and
    magnetic fields in another frame moving relative to the first.\ In the Aether model,
    a stationary charge has a spherically symmetric wave field ($\mathbf{E}$ only, no
    time-averaged angular structure).\ In a boosted frame, the wavefronts are Lorentz-
    contracted and phase-tilted, so the 4-wavevector of the emitted wave acquires the
    transverse components that manifest as~$\mathbf{B}$.\ The frame-dependent mixing of
    flow vector and flow angle that relativity describes as $\mathbf{E}$--$\mathbf{B}$
    mixing is therefore a kinematic consequence of the 4-vector transformation of the
    emitted wave field.

  \item \textbf{Force proportional to velocity:} The magnetic force on a moving charge
    ($\mathbf{F} = q\mathbf{v} \times \mathbf{B}$) depends on the charge's velocity.\ In the
    Aether model, the pressure differential created by the interaction between two wave
    patterns depends on the relative Doppler-anisotropy of their sources.\ A stationary
    charge in a magnetic field experiences no net force because its wave field has no
    Doppler tilt to couple to the external angular structure.\ A moving charge's
    Doppler-anisotropic wave field creates an asymmetric pressure interaction with the
    external vorticity, producing the Lorentz force.
\end{itemize}


\subsection{Electromagnetic Induction}

Faraday's law of electromagnetic induction---a changing magnetic field produces an
electric field, and a changing electric field produces a magnetic field---is the basis
for generators, transformers, and electromagnetic wave propagation.\ The Aether model
provides a physical mechanism for both directions of induction:

\begin{itemize}[leftmargin=2em]
  \item \textbf{Changing $\mathbf{B}$ creates $\mathbf{E}$:} A changing magnetic field is a
    time-varying vorticity envelope in the Aether flow.\ As this rotational pattern
    changes, it buffets the winding pumps of bound rings in its neighborhood and
    drives net coherent wave-field amplitude---which is the electric field.\ The
    changing flow angle displaces Aether in a pattern that produces a net flow vector,
    inducing an electric field in the surrounding region.

  \item \textbf{Changing $\mathbf{E}$ creates $\mathbf{B}$:} A changing electric field is a
    time-varying coherent wave-field amplitude.\ Any time-varying wave source is
    equivalent (via the source-advection mass current) to a moving charge, whose
    Doppler-anisotropic wavefronts carry azimuthal vorticity---producing a magnetic
    field.\ This is Maxwell's displacement current, given a physical interpretation:
    a changing flow vector develops the same angular character as a moving source.
\end{itemize}


\subsection{Electromagnetic Waves}

Under the winding-pump picture, the electric field is already a coherent
spin-sector wave field at every instant, so free-space electromagnetic waves emerge as a
direct consequence: a photon is a propagating excitation of the same spin-wave
mode, whose displacement (flow vector) and vorticity (flow angle) components are
coupled through the medium's wave equation.\ The two aspects regenerate each other as
the wave propagates---an oscillating flow vector develops azimuthal vorticity
(Doppler-anisotropy of the traveling wave), which in turn drives a net transverse
flow, which develops vorticity, and so on---at the speed of light.

This is an oscillation in the Aether medium between two aspects of the same coherent
wave field---its displacement amplitude/direction and its rotational geometry.\ The
electromagnetic wave does not require a separate mechanism from the other Aether
phenomena; it is simply the free-space propagation of the same wave that a stationary
charge already carries in bound, near-field form.

The polarization of electromagnetic waves corresponds to the orientation of the
flow-vector oscillation.\ The fact that electromagnetic waves are transverse (oscillating
perpendicular to the direction of propagation) is consistent with the Aether model:
the flow vector and flow angle disturbances that constitute the wave are perpendicular
to each other and to the direction of propagation, just as observed.\ At the microscopic
level, these wave modes are spin-wave excitations of the Aether's spinor condensate
(Paper~7~\cite{otto-quantum}); through spin--orbit coupling in the condensate
\cite{volovik2003}, they produce effective mass currents that manifest as the velocity
and vorticity fields described here.\ The ring's winding pump is the bound analog
of exactly this spin-wave excitation---a localized, sourced form of the same
fundamental degree of freedom that propagates as free radiation.


% ============================================================

% ============================================================
\section{Summary of Theoretical Properties}
% ============================================================

\subsection{Properties of Magnetism}

\begin{itemize}[leftmargin=2em]
  \item The magnetic field is the angular structure (flow angle) of the Aether velocity
    field, arising from pressure differentials between interacting electron spin vortices.\
    This is distinct from gravity's binding mechanism.\ Magnetism is how the rotational
    geometry of two flows affects the Aether between them; gravity is how each vortex
    concentrates and cycles Aether independently.
  \item Parallel spins repel (opposing flows create high pressure between them); antiparallel
    spins attract (cooperative flows create low pressure between them).
  \item Field lines form closed loops as a topological consequence of vortex circulation.\
    Magnetic monopoles are forbidden by a mathematical identity: $\nabla \cdot (\nabla
    \times \mathbf{v}) \equiv 0$.
  \item Dipole field falls off as $1/r^3$ (partial cancellation of opposing poles at distance),
    compared to gravity's monopole $1/r^2$ falloff.
  \item Moving charges produce magnetic fields because Doppler-anisotropy of the emitted
    spin-phase wave field gives a time-averaged vorticity along the axis of
    motion---kinematic tilting of wavefronts converts flow vector into flow angle
    (Amp\`ere's law).
  \item Ferromagnetic alignment occurs when parallel-spin vortex topology forces electrons
    into antisymmetric spatial configurations, reducing their Coulomb repulsion.\ The
    exchange energy (${\sim}1$~eV) is primarily a Coulomb effect mediated by
    Pauli exclusion (antisymmetric spatial wavefunctions reduce electron--electron
    repulsion), not an amplified magnetic dipole interaction.
  \item Curie temperature marks the transition where thermal Aether turbulence overwhelms
    vortex coupling, destroying domain alignment---analogous to the laminar--turbulent
    transition in classical fluids.
  \item Paramagnetic response reflects partial, statistical alignment of individual vortices
    against thermal disruption; diamagnetic response reflects the back-pressure from
    perturbation of balanced electron-pair flows.
  \item Superconducting flux expulsion (Meissner effect) occurs when Cooper pairs cancel all
    net Aether flow, creating a zero-flow interior that excludes external magnetic fields.
  \item The electron's magnetic moment is antiparallel to its spin, consistent with the
    vortex model's prediction that the circulation pattern (moment) opposes the rotation
    direction (spin angular momentum).
  \item The anomalous magnetic moment ($g$-factor deviation from 2) arises from transient binding
    events modifying the effective vortex flow pattern---connecting magnetism quantitatively
    to the same virtual particle mechanism that produces gravity.
\end{itemize}


\subsection{Properties of Electromagnetism}

\begin{itemize}[leftmargin=2em]
  \item The electric field is the Aether flow vector: the instantaneous direction and
    magnitude of the effective displacement produced by the coherent spin-phase field
    of the bound core ring.\ Charge sign is determined by the rotation coupling between
    ring and surrounding vortex (co-rotation = electron, counter-rotation = positron).
  \item The magnetic field is the flow angle: the rotational geometry (vorticity) of the
    Aether velocity field.\ $\nabla \cdot \mathbf{B} = 0$ is a mathematical identity.
  \item The Coulomb force is the interaction of two spin-phase flux sources: like
    pump directions repel and opposite directions attract ($F \propto q_1 q_2$, the
    flux-source sign rule), with field lines as steady spin-current streamlines and
    Gauss's law as pump conservation.
  \item The electric force falls off as $1/r^2$ (Gauss's law for the conserved
    spin-phase flux; the locked response preserves the static law at all
    separations).
  \item The enormous strength of the electromagnetic force relative to gravity reflects
    a channel asymmetry (the compact spin phase can be pumped without bound, while
    density sources are content-limited) plus quantitative differences between the
    spin-phase pump field (Coulomb) and the bind--unbind residual wave (gravity):
    strength, emission asymmetry, and
    whether the receiver can resonantly respond (two-way vs.\ one-way coupling).
  \item Each electron produces three coexisting Aether effects: coherent spin-phase
    pump field (electric field), flow angle (magnetic moment), and transient binding
    (gravity).
  \item Magnetism from current arises when the flow vector (electric field) of a moving
    charge acquires Doppler-anisotropy, giving nonzero time-averaged vorticity along the
    axis of motion---unifying $\mathbf{E}$ and $\mathbf{B}$ as two aspects of the same
    coherent wave field.
  \item Electromagnetic induction occurs because changing flow angle (changing $\mathbf{B}$)
    drives a net flow vector (creates $\mathbf{E}$), and changing flow vector (changing
    $\mathbf{E}$) develops angular structure (creates $\mathbf{B}$).
  \item Electromagnetic waves are free-space propagations of the same coherent
    spin-wave mode that a stationary charge carries in bound, sourced
    form, with displacement and vorticity components regenerating each other at the
    speed of light.
\end{itemize}


\subsection{Properties of Force Constants}

\begin{itemize}[leftmargin=2em]
  \item Only two independent electromagnetic constants exist at the Aether level: the speed of
    light $c$ (wave propagation speed) and the Coulomb constant $k$ (spin-phase
    flux-source coupling).
  \item The magnetic permeability $\mu_0$ is derived: $\mu_0 = 4\pi k/c^2$.\ It relates the
    flow-vector coupling to the flow-angle coupling via the medium's wave speed.
  \item The magnetic force between moving charges equals the electric force scaled by
    $v_1 v_2/c^2$, where $v_1$ and $v_2$ are the charge velocities---recovered from
    two first-order Doppler factors, one per charge.
  \item The gravitational constant $G$ introduces one additional parameter beyond the Aether
    properties in $k$: $\epsilon$, the fractional asymmetry of the transient bind--unbind cycle.
  \item The ratio $k/G$ decomposes into the product of four independent factors---wave-amplitude ratio, frequency ratio, emission-asymmetry fraction, and receiver-response threshold---and should be derivable from Aether properties (Section~\ref{sec:em-vs-gravity}).
  \item The factor $v^2/c^2$ has a quadruple appearance: time dilation ($\gamma = 1/\sqrt{1 - v^2/c^2}$),
    the magnetic-to-electric force ratio ($F_{\mathrm{mag}}/F_{\mathrm{elec}} = v_1 v_2/c^2$), relativistic
    energy ($E = \gamma mc^2$), and gravitomagnetism ($B_g = g \times v/c^2$, confirmed by GP-B).\ All four arise from the same Aether mechanics: $v^2/c^2$ is
    the ratio of kinetic energy density to the medium's elastic capacity.
  \item The Aether is experimentally confirmed to be perfectly linear (Hookean) in its
    pressure--density response, confirmed by velocity time dilation (to $10^{-8}$ precision
    across velocities from $10^{-6}\,c$ to $0.9994c$) and independently by the Fresnel drag
    coefficient (exact $1 - 1/n^2$ across materials with density ratios up to nearly 6:1).\
    This compressive linearity governs wave propagation; the electromagnetic
    forces in this paper arise from Bernoulli pressure in the
    incompressible-flow limit.
  \item Gravitational time dilation ($\Delta t/t = GM/(rc^2)$) directly connects $G$ to the
    medium: the Aether inflow velocity $\vin = \sqrt{2GM/r}$ produces
    $\vin^2/(2c^2) = GM/(rc^2)$, the time dilation of a static observer in the
    flowing medium.
  \item The fine structure constant ($\alpha \approx 1/137$) measures the coupling
    strength of a single electron's spin-phase pump to the medium; Paper~7(a)
    \cite{otto-boundcore} obtains the closed form
    $\alpha^{-1} = 4\pi[\ln(8\sqrt{2}\,m_A/m_e) - \tfrac12](1 + f_{\mathrm{core}} - 0.017)$.
  \item The gravitational coupling constant ($\alpha_G \approx 1.75 \times 10^{-45}$)
    measures the strength of a single electron's residual binding pressure relative to
    the same quantum scale.
  \item The Planck mass ($m_P \approx 2.18 \times 10^{-8}\;\mathrm{kg}$) represents the mass scale
    at which gravitational binding pressure becomes comparable to quantum-scale Aether
    disturbances.
\end{itemize}



% ============================================================
\section{Experimental Support and Connections to Established Physics}
% ============================================================

The following experimentally confirmed phenomena support or align with the aspects of the theory presented in this paper:

\begin{enumerate}[leftmargin=2em]
  \item \textbf{The anomalous magnetic moment of the electron.} Virtual particle corrections to the
    electron's magnetic moment have been confirmed to approximately thirteen significant digits \cite{fan2023},
    establishing that transient particle--antiparticle formation and annihilation is a real,
    measurable process occurring continuously around charged particles.\ In this theory,
    the anomalous moment arises from transient binding events (the same mechanism that
    produces gravity) modifying the effective Aether vortex flow around the electron---connecting the magnetic and gravitational mechanisms through a single, precisely
    measured quantity.

  \item \textbf{Absence of magnetic monopoles.} Despite extensive experimental searches, no magnetic
    monopole has ever been observed.\ This is a direct consequence of the vortex model:
    circulatory Aether flow inherently produces closed field lines with two poles.
    A monopole would require a topologically impossible radial-only flow from a vortex.

  \item \textbf{The Meissner effect.} Superconductors expel magnetic fields from their interiors.\ In
    this theory, Cooper pairs (opposite-spin electron pairs) cancel all net Aether flow,
    creating a zero-flow region that external magnetic fields cannot penetrate.\ The
    observed exponential decay of the field at the surface (London penetration depth) is
    consistent with the gradual attenuation of external Aether flow by the Cooper pair
    condensate.

  \item \textbf{The $\mu_0 = 4\pi k/c^2$ relationship.} Maxwell's equations \cite{maxwell1865} establish that the speed of
    light equals $1/\sqrt{\epsilon_0 \mu_0}$, linking the electric and magnetic constants
    through $c$.\ In the Aether model, this relationship is not a mathematical coincidence
    but a physical necessity: magnetism is moving electric field, so the magnetic constant
    must equal the electric constant divided by the square of the propagation speed.

\end{enumerate}



% ============================================================
\section{Open Questions and Testable Predictions}
\label{sec:em-open-questions}
% ============================================================

The following areas relevant to this paper remain underdeveloped and represent opportunities for further refinement or falsification:

\begin{enumerate}[leftmargin=2em]
  \item \textbf{Quantitative reproduction of the $g$-factor anomaly.} The electron's anomalous magnetic
    moment ($g = 2.00231930436118$) is known to extraordinary precision and is calculated
    in QED by summing virtual particle loop corrections.\ If the Aether vortex model can
    reproduce this correction from first principles---by calculating how transient binding
    events modify the effective vortex flow pattern---it would provide one of the most
    stringent quantitative tests of the theory.

  \item \textbf{Closing the winding pump.}\
    The concrete ring topology of Paper~7(a)~\cite{otto-boundcore} supplies the
    source through the spin-phase winding pump, normalized by one winding
    quantum per cycle ($q_0 = 2\pi r_{\mathrm{loop}}$), and obtains
    $k\,e^2 = \rhoa\,\kappa_e^2\,r_{\mathrm{loop}}^2/8\pi$---within $0.2\%$ at
    the reference parameter point.\ Remaining first-principles work, developed
    there: the boundary-layer conversion calculation (pre-registered answer:
    duty $2/\pi$ of the layer population per poloidal turn), the
    species-inheritance mechanism for the charge quantum, universal
    synchronization from injection locking, and the isolated-charge energetics
    of the pump field.

  \item \textbf{The fine structure constant.}\ Paper~7(a) reduces $\alpha$ to a
    closed form,
    $\alpha^{-1} = 4\pi[\ln(8\sqrt{2}\,m_A/m_e) - \tfrac12](1+f_{\mathrm{core}}-0.017)$,
    in which $\rhoa$ cancels and $m_A$ enters only logarithmically.\ Exact
    agreement requires $f_{\mathrm{core}} \approx 0.18$--$0.26$, so the GPE
    core-energy calculation is simultaneously the $\alpha$ calculation.

\end{enumerate}



% ============================================================
\section{Mathematical Summary}
\label{sec:em-math-summary}
% ============================================================

This section collects the key equations relevant to the topics covered in this paper.
\noindent\textit{For the complete mathematical summary, see Paper~1: General Overview \cite{otto-overview}.}

\subsection{Fundamental Aether Parameters}
% --------------------------------------------------

The theory posits four microscopic parameters from which all macroscopic
phenomena derive:
%
\begin{center}
\begin{tabular}{@{}llll@{}}
\hline
\textbf{Symbol} & \textbf{Name} & \textbf{Constrained range} & \textbf{Primary constraint} \\
\hline
$\rhoa$  & Ambient density     & $10^{11}$--$10^{12}$~kg/m$^3$         & Electron vortex energy \\
$\Ka$    & Bulk modulus         & $10^{28}$--$10^{29}$~J/m$^3$          & $\Ka = \rhoa\,c^2$ \\
$m_A$    & Aether quantum mass  & $500$--$850$~MeV/$c^2$                & Flux tube width / glueball mass \\
$\xi$    & Healing length       & $0.17$--$0.28$~fm                     & Bridge diameter (lattice QCD) \\
\hline
\end{tabular}
\end{center}

\noindent\textit{The Coulomb constraint of Paper~7(a)~\cite{otto-boundcore}
collapses the $(\rhoa, m_A)$ plane to the curve
$\rhoa/m_A^2 = 8 m_e^2 c^2 e^2 k/(\pi\hbar^4) \approx
3.54\times10^{65}~\mathrm{kg^{-1}\,m^{-3}}$ (e.g.\
$\rhoa \approx 5.0\times10^{11}$~kg/m$^3$ at $m_A = 666$~MeV/$c^2$); the
electron-mass convergence zone is $m_A \approx 616$--$644$~MeV/$c^2$ at the
reference density, with final pinning awaiting the core-energy calculation.}

These four are not independent.\ Two defining relations connect them:
%
\begin{align}
  c &= \sqrt{\Ka / \rhoa}
    \label{eq:sum-speed-of-light} \\
  \xi &= \frac{\hbar}{m_A\,c\,\sqrt{2}}
    \label{eq:sum-healing-length}
\end{align}
%
leaving two free Aether parameters (e.g., $\rhoa$ and $m_A$) from which all
others follow.\ (A third parameter, $\epsilon$, enters through gravity; see below.)


% --------------------------------------------------
\subsection{Equation of State}
% --------------------------------------------------

The Aether obeys a logarithmic equation of state, the unique form that produces
a density-independent bulk modulus (Paper~7~\cite{otto-quantum}):
%
\begin{equation}
  P = \Ka \ln\!\Bigl(\frac{\rho}{\rho_0}\Bigr) + P_0
  \label{eq:sum-eos}
\end{equation}
%
Key consequence: the speed of sound $c_s = \sqrt{\Ka/\rho}$ varies with
density.\ At ambient density, $c_s = c$.\ Note: in gravitational fields,
Bernoulli's equation for free-fall inflow gives $\Delta\rho/\rhoa = 0$
exactly, so $c$ does not vary near gravitating masses; the density
dependence is relevant for material media and for the two-fluid
cosmological structure.\ The electromagnetic forces described in this paper
arise from Bernoulli (dynamic) pressure---the redistribution of kinetic and
static energy in flowing Aether---not from this compressive equation of
state.\ The equation of state governs wave propagation and cosmological
structure; electromagnetic mechanisms require only the incompressible-flow
limit.\ This equation of state guarantees that $\Ka$ is
constant under compression, which is experimentally confirmed by:
%
\begin{itemize}[leftmargin=2em]
  \item Velocity time dilation measurements spanning $10^{-6}\,c$ to $0.9994\,c$
    (precision $\sim 10^{-8}$).
  \item The Fresnel drag coefficient $1 - 1/n^2$, exact across materials with
    density ratios up to $\sim$6:1.
\end{itemize}


% --------------------------------------------------
\subsection{Electromagnetic Equations}
% --------------------------------------------------

\paragraph{Electric field (flow vector).}
The electric field is the Aether flow vector---the instantaneous direction and
magnitude of the effective displacement produced by the coherent spin-phase field
of the bound core ring (Section~\ref{sec:electric-field}).\ The Coulomb force between
charges $q_1, q_2$ at separation $r$:
%
\begin{equation}
  F_{\mathrm{elec}} = k\,\frac{q_1 q_2}{r^2}
  \label{eq:sum-coulomb}
\end{equation}
%
with $k$ derived in Paper~7(a)~\cite{otto-boundcore} from the ring's winding
pump: $k\,e^2 = \rhoa\,\kappa_e^2\,r_{\mathrm{loop}}^2/8\pi$, equivalently
$\alpha^{-1} = 4\pi[\ln(8\sqrt{2}\,m_A/m_e) - \tfrac12](1+f_{\mathrm{core}}-0.017)$.

\paragraph{Magnetic force (flow angle).}
Magnetism arises from the angular structure of the Aether flow---the rotational
geometry of the velocity field.\ The magnetic force between charges moving at
velocities $v_1, v_2$:
%
\begin{equation}
  F_{\mathrm{mag}} = F_{\mathrm{elec}}\,\frac{v_1\,v_2}{c^2}
  \label{eq:sum-mag-force}
\end{equation}

\paragraph{Permeability.}
The magnetic permeability relates the spin-phase (flow-vector) coupling
to the flow-angle coupling via the medium's wave speed, and is a derived quantity, not
a fundamental constant:
%
\begin{equation}
  \mu_0 = \frac{4\pi\,k}{c^2}
  \label{eq:sum-mu0}
\end{equation}


% --------------------------------------------------

\section*{Acknowledgments}
\addcontentsline{toc}{section}{Acknowledgments}

The author acknowledges the use of Claude (Anthropic) for \LaTeX{} formatting
and editorial assistance.

\bigskip


% ============================================================
% BIBLIOGRAPHY
% ============================================================

\begin{thebibliography}{99}

\bibitem{maxwell1865}
Maxwell, J.~C.,
``A Dynamical Theory of the Electromagnetic Field,''
\textit{Phil.\ Trans.\ Royal Society of London}, \textbf{155}, 459--512 (1865).


\bibitem{curie1895}
Curie, P.,
``Propri\'{e}t\'{e}s magn\'{e}tiques des corps \`{a} diverses temp\'{e}ratures,''
\textit{Annales de chimie et de physique}, 7th series, \textbf{5}, 289--405 (1895).




\bibitem{fan2023}
Fan, X., Myers, T.~G., Sukra, B.~A.~D., and Gabrielse, G.,
``Measurement of the Electron Magnetic Moment,''
\textit{Physical Review Letters}, \textbf{130}, 071801 (2023).

\bibitem{otto-overview}
Otto, E.,
``Unifying Theory of Dynamic Aether Mechanics: General Overview,''
(2026).

\bibitem{otto-cosmology}
Otto, E.,
``Unifying Theory of Dynamic Aether Mechanics: Cosmology,''
(2026).

\bibitem{otto-quantum}
Otto, E.,
``Unifying Theory of Dynamic Aether Mechanics: Quantum Mechanics,''
(2026).

\bibitem{otto-boundcore}
Otto, E.,
``Unifying Theory of Dynamic Aether Mechanics: Bound Core Electron Model,''
(2026).

\bibitem{otto-gravity}
Otto, E.,
``Unifying Theory of Dynamic Aether Mechanics: Gravity,''
(2026).

\bibitem{bjerknes1906}
Bjerknes, V.~F.~K.,
\textit{Fields of Force}
(Columbia University Press, New York, 1906).\
See also: T.~G.~Leighton, \textit{The Acoustic Bubble}
(Academic Press, London, 1994), Chapter~4, for the modern treatment of
secondary Bjerknes forces between phase-locked pulsating sources in a fluid medium.

\bibitem{volovik2003}
Volovik, G.~E.,
\textit{The Universe in a Helium Droplet}
(Oxford University Press, Oxford, 2003).
\end{thebibliography}

\end{document}